\documentclass[a4paper,10pt,BCOR10mm,oneside,headsepline]{scrartcl}
\usepackage[portuges]{babel}
\usepackage[utf8]{inputenc}
\usepackage{amsmath}
\usepackage{wasysym}% provides \ocircle and \Box
\usepackage{enumitem}% easy control of topsep and leftmargin for lists
\usepackage{color}% used for background color
\usepackage{forloop}% used for \Qrating and \Qlines
\usepackage{ifthen}% used for \Qitem and \QItem
\usepackage{typearea}
\usepackage{graphicx}
\usepackage{booktabs}
\usepackage{listings}
\usepackage{color}
\usepackage{algorithm}
\usepackage{algorithmic}
\usepackage{amssymb}
\usepackage{minted}
\usepackage{mdframed}
\usepackage{wasysym}
\usepackage{xcolor}
\usepackage{fontspec}
%\usepackage{tgtermes}

\areaset{17cm}{26cm}
\setlength{\topmargin}{-2cm}

\definecolor{codegreen}{rgb}{0,0.6,0}
\definecolor{codegray}{rgb}{0.5,0.5,0.5}
\definecolor{codepurple}{rgb}{0.58,0,0.82}
\definecolor{backcolour}{rgb}{0.95,0.95,0.92}
 
\lstdefinestyle{mystyle}{
    backgroundcolor=\color{backcolour},   
    commentstyle=\color{codegreen},
    keywordstyle=\color{magenta},
    numberstyle=\tiny\color{codegray},
    stringstyle=\color{codepurple},
    basicstyle=\footnotesize,
    breakatwhitespace=false,         
    breaklines=true,                 
    captionpos=b,                    
    keepspaces=true,                 
    numbers=left,                    
    numbersep=5pt,                  
    showspaces=false,                
    showstringspaces=false,
    showtabs=false,                  
    tabsize=2
}
 
\lstset{style=mystyle}

\mdfsetup{
    skipabove=12pt,
    skipbelow=12pt,
    innertopmargin=8pt,
    innerbottommargin=8pt,
    innerleftmargin=10pt,
    innerrightmargin=10pt,
    splittopskip=12pt,
    splitbottomskip=12pt,
}

\algsetup{linenodelimiter=.,linenosize=\tiny}

\renewcommand{\algorithmicrequire}{\textbf{Entrada:}}
\renewcommand{\algorithmicensure}{\textbf{Saída:}}
\floatname{algorithm}{Algoritmo}

\setmainfont{TeX Gyre Termes}
\setmonofont{APL 385 Unicode}[Scale=0.9]

%% \Qq = Questionaire question. Oh, this is just too simple. It helps
%% making it easy to globally change the appearance of questions.
\newcommand{\Qq}[1]{\textbf{#1}}
%\newcommand{\Qq}[1]{{#1}}

%% \QO = Circle or box to be ticked. Used both by direct call and by
%% \Qrating and \Qlist.
\newcommand{\QO}{$\Box$}% or: $\ocircle$

%% \Qrating = Automatically create a rating scale with NUM steps, like
%% this: 0--0--0--0--0.
\newcounter{qr}
\newcommand{\Qrating}[1]{\QO\forloop{qr}{1}{\value{qr} < #1}{---\QO}}

%% \Qline = Again, this is very simple. It helps setting the line
%% thickness globally. Used both by direct call and by \Qlines.
\newcommand{\Qline}[1]{\noindent\rule{#1}{0.6pt}}

%% \Qlines = Insert NUM lines with width=\linewith. You can change the
%% \vskip value to adjust the spacing.
\newcounter{ql}
\newcommand{\Qlines}[1]{\forloop{ql}{0}{\value{ql}<#1}{\vskip0em\Qline{\linewidth}}}

%% \Qlist = This is an environment very similar to itemize but with
%% \QO in front of each list item. Useful for classical multiple
%% choice. Change leftmargin and topsep accourding to your taste.
\newenvironment{Qlist}{%
\renewcommand{\labelitemi}{\QO}
\begin{itemize}[leftmargin=1.5em,topsep=-.5em]
}{%
\end{itemize}
}

%% \Qtab = A "tabulator simulation". The first argument is the
%% distance from the left margin. The second argument is content which
%% is indented within the current row.
\newlength{\qt}
\newcommand{\Qtab}[2]{
\setlength{\qt}{\linewidth}
\addtolength{\qt}{-#1}
\hfill\parbox[t]{\qt}{\raggedright #2}
}

%% \Qitem = Item with automatic numbering. The first optional argument
%% can be used to create sub-items like 2a, 2b, 2c, ... The item
%% number is increased if the first argument is omitted or equals 'a'.
%% You will have to adjust this if you prefer a different numbering
%% scheme. Adjust topsep and leftmargin as needed.
\newcounter{itemnummer}
\newcommand{\Qitem}[2][]{% #1 optional, #2 notwendig
\ifthenelse{\equal{#1}{}}{\stepcounter{itemnummer}}{}
\ifthenelse{\equal{#1}{a}}{\stepcounter{itemnummer}}{}
\begin{enumerate}[topsep=2pt,leftmargin=2.8em]
\item[\textbf{\arabic{itemnummer}#1.}] #2
\end{enumerate}
}

%% \QItem = Like \Qitem but with alternating background color. This
%% might be error prone as I hard-coded some lengths (-5.25pt and
%% -3pt)! I do not yet understand why I need them.
\definecolor{bgodd}{rgb}{0.8,0.8,0.8}
\definecolor{bgeven}{rgb}{0.9,0.9,0.9}
\newcounter{itemoddeven}
\newlength{\gb}
\newcommand{\QItem}[2][]{% #1 optional, #2 notwendig
\setlength{\gb}{\linewidth}
\addtolength{\gb}{-5.25pt}
\ifthenelse{\equal{\value{itemoddeven}}{0}}{%
\noindent\colorbox{bgeven}{\hskip-3pt\begin{minipage}{\gb}\Qitem[#1]{#2}\end{minipage}}%
\stepcounter{itemoddeven}%
}{%
\noindent\colorbox{bgodd}{\hskip-3pt\begin{minipage}{\gb}\Qitem[#1]{#2}\end{minipage}}%
\setcounter{itemoddeven}{0}%
}
}

\newcommand{\rawcode}[1]{\texttt{#1}}
\newcommand{\code}[2]{\mintinline{#1}{#2}}
\newcommand{\file}[1]{\texttt{\underline{#1}}}

\renewcommand{\theFancyVerbLine}{\tiny \texttt{\textcolor{gray}{\arabic{FancyVerbLine}}}}

\newcommand{\inputcode}[2]{\begin{mdframed}\inputminted[fontsize=\footnotesize,linenos,numbersep=4pt]{#1}{#2}\end{mdframed}}

\newenvironment{solution}{%
    \begin{mdframed}%
        \noindent\textbf{Solu\c{c}\~ao:}\par\medskip
}{%
    \end{mdframed}%
}
\newcommand{\inputsnippet}[4]{\begin{mdframed}[shadow=true,shadowsize=3pt,shadowcolor=gray]\inputminted[fontsize=\footnotesize,linenos,firstline=#2,lastline=#3,numbersep=2pt]{#1}{#4}\end{mdframed}}

\newcommand{\inputsyntax}[2]{\inputminted[fontsize=\footnotesize]{#1}{#2}}


\mdfsetup{nobreak=true}

\usepackage{tikz}
\usetikzlibrary{arrows.meta}
\tikzset{
    tmstate/.style={draw,circle,inner sep=1pt,minimum size=6.5mm,font=\footnotesize},
    abnode/.style={draw,circle,inner sep=1pt,minimum size=7mm,font=\footnotesize}
}

\begin{document}

\begin{center}
	\includegraphics[scale=0.45]{fga}
	% fga.ps: 1087x103 pixel, 72dpi, 38.35x3.63 cm, bb=0 0 1087 103
\end{center}


\vspace{0.1in}

\begin{center}
    \textbf{\huge Paradigmas de Programação}

    \vspace{0.1in}
    \textbf{\Large Máquinas de Turing: Lista 2}
\end{center}\vskip1em

\renewcommand{\Qq}[1]{{#1}}

\section*{Parte A - Definição de Máquinas de Turing}

\begin{mdframed}
\noindent\textbf{Convenções:}\par\medskip
Nos fluxogramas, cada seta tem rótulo $s:a$, com $s\in\{0,1\}$ ($1$ = traço, $0$ = branco) e
$a\in\{L,R,B,1\}$ (mover à esquerda, mover à direita, apagar ou escrever traço). A máquina
inicia em $q_1$ e para quando não há transição para o par (estado, símbolo); estados de parada
aparecem com círculo duplo.
\end{mdframed}

\Qitem{ \Qq{Considere uma fita contendo um bloco de $n$ traços, seguido por um espaço, seguido por
        um bloco de $m$ traços, seguido por um espaço, seguido por um bloco de $k$ traços, sendo
        que o resto da fita está em branco. Especifique uma máquina de Turing, por meio de um
        fluxograma, que, quando iniciada no traço mais à esquerda, finalmente para, sem ter
        escrito nem apagado nada...}

    \vspace{0.1in}
    \begin{Qlist}
        \item[(a)] ...no traço mais à esquerda do segundo bloco.

        \begin{solution}
        $q_1$ percorre o 1º bloco e avança sobre o espaço até o primeiro traço do 2º bloco, parando
        em $q_2$.

        \begin{center}
        \begin{tikzpicture}[>=Stealth,node distance=1cm,every node/.style={}]
  \node[tmstate] (q1) at (0.00,0.00) {$q_{1}$};
  \node[tmstate,double] (q2) at (2.40,0.00) {$q_{2}$};
  \draw[->] (q1) edge[loop above] node[above] {\scriptsize $1:R$} (q1);
  \draw[->] (q1) edge node[above,midway] {\scriptsize $0:R$} (q2);
\end{tikzpicture}
        \end{center}
        \end{solution}

        \item[(b)] ...no traço mais à esquerda do terceiro bloco.

        \begin{solution}
        $q_1$ e $q_2$ atravessam os dois primeiros blocos e seus espaços seguintes, parando em $q_3$
        sobre o primeiro traço do 3º bloco.

        \begin{center}
        \begin{tikzpicture}[>=Stealth,node distance=1cm,every node/.style={}]
  \node[tmstate] (q1) at (0.00,0.00) {$q_{1}$};
  \node[tmstate] (q2) at (2.40,0.00) {$q_{2}$};
  \node[tmstate,double] (q3) at (4.80,0.00) {$q_{3}$};
  \draw[->] (q1) edge[loop above] node[above] {\scriptsize $1:R$} (q1);
  \draw[->] (q1) edge node[above,midway] {\scriptsize $0:R$} (q2);
  \draw[->] (q2) edge[loop above] node[above] {\scriptsize $1:R$} (q2);
  \draw[->] (q2) edge node[above,midway] {\scriptsize $0:R$} (q3);
\end{tikzpicture}
        \end{center}
        \end{solution}

        \item[(c)] ...no traço mais à direita do segundo bloco.

        \begin{solution}
        $q_1$ e $q_2$ percorrem os dois primeiros blocos; ao ler o espaço após o 2º bloco, $q_2$
        recua um passo para o último traço desse bloco ($q_3$).

        \begin{center}
        \begin{tikzpicture}[>=Stealth,node distance=1cm,every node/.style={}]
  \node[tmstate] (q1) at (0.00,0.00) {$q_{1}$};
  \node[tmstate] (q2) at (2.40,0.00) {$q_{2}$};
  \node[tmstate,double] (q3) at (4.80,0.00) {$q_{3}$};
  \draw[->] (q1) edge[loop above] node[above] {\scriptsize $1:R$} (q1);
  \draw[->] (q1) edge node[above,midway] {\scriptsize $0:R$} (q2);
  \draw[->] (q2) edge[loop above] node[above] {\scriptsize $1:R$} (q2);
  \draw[->] (q2) edge node[above,midway] {\scriptsize $0:L$} (q3);
\end{tikzpicture}
        \end{center}
        \end{solution}

        \item[(d)] ...no traço mais à direita do terceiro bloco.

        \begin{solution}
        $q_1$, $q_2$ e $q_3$ percorrem os três blocos; ao ler o branco após o 3º bloco, $q_3$ recua
        um passo para o último traço desse bloco ($q_4$).

        \begin{center}
        \begin{tikzpicture}[>=Stealth,node distance=1cm,every node/.style={}]
  \node[tmstate] (q1) at (0.00,0.00) {$q_{1}$};
  \node[tmstate] (q2) at (2.40,0.00) {$q_{2}$};
  \node[tmstate] (q3) at (4.80,0.00) {$q_{3}$};
  \node[tmstate,double] (q4) at (7.20,0.00) {$q_{4}$};
  \draw[->] (q1) edge[loop above] node[above] {\scriptsize $1:R$} (q1);
  \draw[->] (q1) edge node[above,midway] {\scriptsize $0:R$} (q2);
  \draw[->] (q2) edge[loop above] node[above] {\scriptsize $1:R$} (q2);
  \draw[->] (q2) edge node[above,midway] {\scriptsize $0:R$} (q3);
  \draw[->] (q3) edge[loop above] node[above] {\scriptsize $1:R$} (q3);
  \draw[->] (q3) edge node[above,midway] {\scriptsize $0:L$} (q4);
\end{tikzpicture}
        \end{center}
        \end{solution}
    \end{Qlist}
}

\Qitem{ \Qq{Especifique, por meio de um fluxograma, uma máquina de Turing que, iniciando com a fita
        descrita na questão anterior, finalmente para no traço mais à esquerda da fita, a qual
        deve agora conter...}

    \vspace{0.1in}
    \begin{Qlist}
        \item[(a)] um bloco de $n$ traços, seguido por um espaço em branco, seguido por um
            bloco de $m + 1$ traços, seguido por um espaço em branco, seguido por um bloco de $k$
            traços.

        \begin{solution}
        Desloca-se o 1º bloco uma posição para a esquerda (escrevendo um traço antes do início e
        apagando seu último traço) e preenche-se o espaço original entre os blocos, acrescentando
        um traço ao 2º bloco. Ao fim, a cabeça volta ao traço mais à esquerda ($q_8$).

        \begin{center}
        \begin{tikzpicture}[>=Stealth,node distance=1cm,every node/.style={}]
  \node[tmstate] (q1) at (0.00,0.00) {$q_{1}$};
  \node[tmstate] (q2) at (2.40,0.00) {$q_{2}$};
  \node[tmstate] (q3) at (4.80,0.00) {$q_{3}$};
  \node[tmstate] (q4) at (7.20,0.00) {$q_{4}$};
  \node[tmstate] (q5) at (7.20,-2.00) {$q_{5}$};
  \node[tmstate] (q6) at (4.80,-2.00) {$q_{6}$};
  \node[tmstate] (q7) at (2.40,-2.00) {$q_{7}$};
  \node[tmstate,double] (q8) at (0.00,-2.00) {$q_{8}$};
  \draw[->] (q1) edge node[above,midway] {\scriptsize $1:L$} (q2);
  \draw[->] (q2) edge node[above,midway] {\scriptsize $0:1$} (q3);
  \draw[->] (q3) edge[loop above] node[above] {\scriptsize $1:R$} (q3);
  \draw[->] (q3) edge node[above,midway] {\scriptsize $0:1$} (q4);
  \draw[->] (q4) edge node[right,midway] {\scriptsize $1:L$} (q5);
  \draw[->] (q5) edge node[above,midway] {\scriptsize $1:B$} (q6);
  \draw[->] (q6) edge node[above,midway] {\scriptsize $0:L$} (q7);
  \draw[->] (q7) edge[loop below] node[below] {\scriptsize $1:L$} (q7);
  \draw[->] (q7) edge node[above,midway] {\scriptsize $0:R$} (q8);
\end{tikzpicture}
        \end{center}
        \end{solution}

        \item[(b)] um bloco de $n$ traços, seguido por um espaço em branco, seguido por um
            bloco de $m - 1$ traços, seguido por um espaço em branco, seguido por um bloco de $k$
            traços.

        \begin{solution}
        Apaga-se o primeiro traço do 1º bloco, preenche-se o espaço original e apaga-se o primeiro
        traço do 2º bloco, deslocando o 1º bloco uma posição para a direita sem deixar espaço extra. Ao fim, volta-se ao traço mais à esquerda ($q_8$).

        \begin{center}
        \begin{tikzpicture}[>=Stealth,node distance=1cm,every node/.style={}]
  \node[tmstate] (q1) at (0.00,0.00) {$q_{1}$};
  \node[tmstate] (q2) at (2.40,0.00) {$q_{2}$};
  \node[tmstate] (q3) at (4.80,0.00) {$q_{3}$};
  \node[tmstate] (q4) at (7.20,0.00) {$q_{4}$};
  \node[tmstate] (q5) at (7.20,-2.00) {$q_{5}$};
  \node[tmstate] (q6) at (4.80,-2.00) {$q_{6}$};
  \node[tmstate] (q7) at (2.40,-2.00) {$q_{7}$};
  \node[tmstate,double] (q8) at (0.00,-2.00) {$q_{8}$};
  \draw[->] (q1) edge node[above,midway] {\scriptsize $1:B$} (q2);
  \draw[->] (q2) edge node[above,midway] {\scriptsize $0:R$} (q3);
  \draw[->] (q3) edge[loop above] node[above] {\scriptsize $1:R$} (q3);
  \draw[->] (q3) edge node[above,midway] {\scriptsize $0:1$} (q4);
  \draw[->] (q4) edge node[right,midway] {\scriptsize $1:R$} (q5);
  \draw[->] (q5) edge node[above,midway] {\scriptsize $1:B$} (q6);
  \draw[->] (q6) edge node[above,midway] {\scriptsize $0:L$} (q7);
  \draw[->] (q7) edge[loop below] node[below] {\scriptsize $1:L$} (q7);
  \draw[->] (q7) edge node[above,midway] {\scriptsize $0:R$} (q8);
\end{tikzpicture}
        \end{center}
        \end{solution}
    \end{Qlist}
}

\Qitem{ \Qq{Especifique, por meio de um fluxograma, uma máquina de Turing que compute a função}

    \begin{solution}
    Pela convenção padrão, a fita inicia com blocos de $x$ e $y$ traços ($x,y\geq 1$) separados
    por um espaço, com a cabeça no traço mais à esquerda, e deve parar no traço mais à esquerda
    do bloco resultado (com o resto em branco).

    Ambas as máquinas contém o mesmo laço de emparelhamento ($q_1$ a $q_{11}$): a cada rodada,
    move-se um traço do 1º bloco para um acumulador $R$ à esquerda e apaga-se um traço do 2º
    bloco. Após $r=\min(x,y)$ rodadas, a fita fica na forma
    \[
        \underbrace{1\cdots 1}_{R:\ r}\ \ 0_b\ \ \underbrace{1\cdots 1}_{X':\ x-r}\ \ 0\ \ 
        \underbrace{1\cdots 1}_{Y':\ y-r},
    \]
    em que $X'$ (se $x\leq y$) ou $Y'$ (se $y<x$) está vazio.
    \end{solution}

    \vspace{0.1in}
    \begin{Qlist}
        \item[(a)] $\min(x, y) = $ o menor de $x$ e $y$.

        \begin{solution}
        Ao esgotar $X'$ ($q_1\to q_{12}$) ou $Y'$ ($q_3\to q_{16}$), apaga-se o bloco restante
        ($q_{12}$--$q_{13}$ ou $q_{16}$--$q_{18}$) e retorna-se ao traço mais à esquerda de $R$
        ($q_{15}\to q_{19}$), que contém $\min(x,y)$ traços.

        \begin{center}
        \resizebox{\linewidth}{!}{\begin{tikzpicture}[>=Stealth,node distance=1cm,every node/.style={}]
  \node[tmstate] (q1) at (0.00,0.00) {$q_{1}$};
  \node[tmstate] (q2) at (2.20,0.00) {$q_{2}$};
  \node[tmstate] (q3) at (4.40,0.00) {$q_{3}$};
  \node[tmstate] (q4) at (6.60,0.00) {$q_{4}$};
  \node[tmstate] (q5) at (8.80,0.00) {$q_{5}$};
  \node[tmstate] (q6) at (8.80,-2.00) {$q_{6}$};
  \node[tmstate] (q7) at (6.60,-2.00) {$q_{7}$};
  \node[tmstate] (q8) at (4.40,-2.00) {$q_{8}$};
  \node[tmstate] (q9) at (2.20,-2.00) {$q_{9}$};
  \node[tmstate] (q10) at (0.00,-2.00) {$q_{10}$};
  \node[tmstate] (q11) at (-2.20,-1.00) {$q_{11}$};
  \node[tmstate] (q12) at (0.00,2.00) {$q_{12}$};
  \node[tmstate] (q13) at (-2.20,2.00) {$q_{13}$};
  \node[tmstate] (q14) at (0.00,4.00) {$q_{14}$};
  \node[tmstate] (q15) at (3.30,4.00) {$q_{15}$};
  \node[tmstate,double] (q19) at (3.30,6.00) {$q_{19}$};
  \node[tmstate] (q16) at (4.40,2.00) {$q_{16}$};
  \node[tmstate] (q17) at (6.60,2.00) {$q_{17}$};
  \node[tmstate] (q18) at (8.80,2.00) {$q_{18}$};
  \draw[->] (q1) edge node[above,midway] {\scriptsize $1:R$} (q2);
  \draw[->] (q1) edge node[right,midway] {\scriptsize $0:R$} (q12);
  \draw[->] (q2) edge[loop above] node[above] {\scriptsize $1:R$} (q2);
  \draw[->] (q2) edge node[above,midway] {\scriptsize $0:R$} (q3);
  \draw[->] (q3) edge node[above,midway] {\scriptsize $1:R$} (q4);
  \draw[->] (q3) edge node[right,midway] {\scriptsize $0:L$} (q16);
  \draw[->] (q4) edge[loop above] node[above] {\scriptsize $1:R$} (q4);
  \draw[->] (q4) edge node[above,midway] {\scriptsize $0:L$} (q5);
  \draw[->] (q5) edge node[right,midway] {\scriptsize $1:B$} (q6);
  \draw[->] (q6) edge node[above,midway] {\scriptsize $0:L$} (q7);
  \draw[->] (q7) edge[loop below] node[below] {\scriptsize $1:L$} (q7);
  \draw[->] (q7) edge node[above,midway] {\scriptsize $0:L$} (q8);
  \draw[->] (q8) edge[loop below] node[below] {\scriptsize $1:L$} (q8);
  \draw[->] (q8) edge node[above,midway] {\scriptsize $0:1$} (q9);
  \draw[->] (q9) edge node[above,midway] {\scriptsize $1:R$} (q10);
  \draw[->] (q10) edge node[above,sloped,midway] {\scriptsize $1:B$} (q11);
  \draw[->] (q11) edge node[above,sloped,midway] {\scriptsize $0:R$} (q1);
  \draw[->] (q12) edge[bend left] node[below] {\scriptsize $1:B$} (q13);
  \draw[->] (q12) edge node[right,midway] {\scriptsize $0:L$} (q14);
  \draw[->] (q13) edge[bend left] node[above] {\scriptsize $0:R$} (q12);
  \draw[->] (q14) edge[loop above] node[above] {\scriptsize $0:L$} (q14);
  \draw[->] (q14) edge node[above,midway] {\scriptsize $1:L$} (q15);
  \draw[->] (q15) edge[loop below] node[below] {\scriptsize $1:L$} (q15);
  \draw[->] (q15) edge node[right,midway] {\scriptsize $0:R$} (q19);
  \draw[->] (q16) edge node[above,midway] {\scriptsize $0:L$} (q17);
  \draw[->] (q17) edge[bend left] node[above] {\scriptsize $1:B$} (q18);
  \draw[->] (q17) edge node[above,sloped,midway] {\scriptsize $0:L$} (q15);
  \draw[->] (q18) edge[bend left] node[below] {\scriptsize $0:L$} (q17);
\end{tikzpicture}}
        \end{center}
        \end{solution}

        \item[(b)] $\max(x, y) = $ o maior de $x$ e $y$.

        \begin{solution}
        Ao esgotar $X'$ ou $Y'$, preenchem-se o espaço separador e $0_b$ ($q_{12},q_{13}$), unindo
        $R$ ao excesso restante em um bloco de $\max(x,y)+2$ traços. Por fim, volta-se ao início
        desse bloco e apagam-se seus dois primeiros traços ($q_{14}$--$q_{19}$).

        \begin{center}
        \resizebox{\linewidth}{!}{\begin{tikzpicture}[>=Stealth,node distance=1cm,every node/.style={}]
  \node[tmstate] (q1) at (0.00,0.00) {$q_{1}$};
  \node[tmstate] (q2) at (2.20,0.00) {$q_{2}$};
  \node[tmstate] (q3) at (4.40,0.00) {$q_{3}$};
  \node[tmstate] (q4) at (6.60,0.00) {$q_{4}$};
  \node[tmstate] (q5) at (8.80,0.00) {$q_{5}$};
  \node[tmstate] (q6) at (8.80,-2.00) {$q_{6}$};
  \node[tmstate] (q7) at (6.60,-2.00) {$q_{7}$};
  \node[tmstate] (q8) at (4.40,-2.00) {$q_{8}$};
  \node[tmstate] (q9) at (2.20,-2.00) {$q_{9}$};
  \node[tmstate] (q10) at (0.00,-2.00) {$q_{10}$};
  \node[tmstate] (q11) at (-2.20,-1.00) {$q_{11}$};
  \node[tmstate] (q12) at (0.00,2.00) {$q_{12}$};
  \node[tmstate] (q13) at (4.40,2.00) {$q_{13}$};
  \node[tmstate] (q14) at (0.00,4.00) {$q_{14}$};
  \node[tmstate] (q15) at (2.20,4.00) {$q_{15}$};
  \node[tmstate] (q16) at (4.40,4.00) {$q_{16}$};
  \node[tmstate] (q17) at (6.60,4.00) {$q_{17}$};
  \node[tmstate] (q18) at (8.80,4.00) {$q_{18}$};
  \node[tmstate,double] (q19) at (11.00,4.00) {$q_{19}$};
  \draw[->] (q1) edge node[above,midway] {\scriptsize $1:R$} (q2);
  \draw[->] (q1) edge node[right,midway] {\scriptsize $0:1$} (q12);
  \draw[->] (q2) edge[loop above] node[above] {\scriptsize $1:R$} (q2);
  \draw[->] (q2) edge node[above,midway] {\scriptsize $0:R$} (q3);
  \draw[->] (q3) edge node[above,midway] {\scriptsize $1:R$} (q4);
  \draw[->] (q3) edge node[right,midway] {\scriptsize $0:L$} (q13);
  \draw[->] (q4) edge[loop above] node[above] {\scriptsize $1:R$} (q4);
  \draw[->] (q4) edge node[above,midway] {\scriptsize $0:L$} (q5);
  \draw[->] (q5) edge node[right,midway] {\scriptsize $1:B$} (q6);
  \draw[->] (q6) edge node[above,midway] {\scriptsize $0:L$} (q7);
  \draw[->] (q7) edge[loop below] node[below] {\scriptsize $1:L$} (q7);
  \draw[->] (q7) edge node[above,midway] {\scriptsize $0:L$} (q8);
  \draw[->] (q8) edge[loop below] node[below] {\scriptsize $1:L$} (q8);
  \draw[->] (q8) edge node[above,midway] {\scriptsize $0:1$} (q9);
  \draw[->] (q9) edge node[above,midway] {\scriptsize $1:R$} (q10);
  \draw[->] (q10) edge node[above,sloped,midway] {\scriptsize $1:B$} (q11);
  \draw[->] (q11) edge node[above,sloped,midway] {\scriptsize $0:R$} (q1);
  \draw[->] (q12) edge[loop left] node[left] {\scriptsize $1:L$} (q12);
  \draw[->] (q12) edge node[right,midway] {\scriptsize $0:1$} (q14);
  \draw[->] (q13) edge node[above,midway] {\scriptsize $0:1$} (q12);
  \draw[->] (q14) edge[loop left] node[left] {\scriptsize $1:L$} (q14);
  \draw[->] (q14) edge node[above,midway] {\scriptsize $0:R$} (q15);
  \draw[->] (q15) edge node[above,midway] {\scriptsize $1:B$} (q16);
  \draw[->] (q16) edge node[above,midway] {\scriptsize $0:R$} (q17);
  \draw[->] (q17) edge node[above,midway] {\scriptsize $1:B$} (q18);
  \draw[->] (q18) edge node[above,midway] {\scriptsize $0:R$} (q19);
\end{tikzpicture}}
        \end{center}
        \end{solution}
    \end{Qlist}
}

\section*{Parte B - Computabilidade}

\Qitem{ \Qq{Identifique o número inteiro associado a máquina de Turing $M_i$. Em seguida, 
        desenhe o fluxograma desta máquina e explique o seu comportamento.}

    \begin{solution}
    Uma máquina com $n$ estados ativos ($q_1,\ldots,q_n$) e parada em $q_{n+1}$ é codificada por
    $2^{a_1}3^{a_2}\cdots p_{4n}^{a_{4n}}$, onde $(a_{2k-1},a_{2k})$ dá (instrução, próximo estado)
    para $q_1S_0, q_1S_1, q_2S_0, q_2S_1,\ldots$, com $S_0,S_1,L,R \mapsto 1,2,3,4$.
    \begin{center}
    \resizebox{\linewidth}{!}{%
    \begin{tabular}{c|ccccccc}
    $i$ & 1 & 2 & 3 & 4 & 5 & 6 & 7 \\ \midrule
    $(a_1,a_2,a_3,a_4)$ & $(1,1,1,1)$ & $(2,1,1,1)$ & $(1,2,1,1)$ & $(3,1,1,1)$ &
        $(1,1,2,1)$ & $(2,2,1,1)$ & $(1,1,1,2)$ \\
    número & $210$ & $420$ & $630$ & $840$ & $1050$ & $1260$ & $1470$
    \end{tabular}}
    \end{center}
    \end{solution}

    \vspace{0.1in}
    \begin{Qlist}
        \item[(a)] $i = 4$

        \begin{solution}
        $840 = 2^3\cdot 3\cdot 5\cdot 7$, isto é, $(a_1,a_2,a_3,a_4)=(3,1,1,1)$:
        $q_1S_0Lq_1$ e $q_1S_1S_0q_1$.
        \begin{center}
        \begin{tikzpicture}[>=Stealth,node distance=1cm,every node/.style={}]
  \node[tmstate] (q1) at (0.00,0.00) {$q_{1}$};
  \draw[->] (q1) edge[loop above] node[above] {\scriptsize $0:L,\ 1:B$} (q1);
\end{tikzpicture}
        \end{center}
        Apaga o traço inicial e move-se indefinidamente para a esquerda sobre brancos; nunca para.
        \end{solution}

        \item[(b)] $i = 5$

        \begin{solution}
        $1050 = 2\cdot 3\cdot 5^2\cdot 7$, isto é, $(a_1,a_2,a_3,a_4)=(1,1,2,1)$:
        $q_1S_0S_0q_1$ e $q_1S_1S_1q_1$.
        \begin{center}
        \begin{tikzpicture}[>=Stealth,node distance=1cm,every node/.style={}]
  \node[tmstate] (q1) at (0.00,0.00) {$q_{1}$};
  \draw[->] (q1) edge[loop above] node[above] {\scriptsize $0:B,\ 1:1$} (q1);
\end{tikzpicture}
        \end{center}
        Reescreve o símbolo examinado sem mover a cabeça nem mudar de estado; nunca para.
        \end{solution}

        \item[(c)] $i = 6$

        \begin{solution}
        $1260 = 2^2\cdot 3^2\cdot 5\cdot 7$, isto é, $(a_1,a_2,a_3,a_4)=(2,2,1,1)$:
        $q_1S_0S_1q_2$ e $q_1S_1S_0q_1$.
        \begin{center}
        \begin{tikzpicture}[>=Stealth,node distance=1cm,every node/.style={}]
  \node[tmstate] (q1) at (0.00,0.00) {$q_{1}$};
  \node[tmstate,double] (q2) at (2.80,0.00) {$q_{2}$};
  \draw[->] (q1) edge node[above,midway] {\scriptsize $0:1$} (q2);
  \draw[->] (q1) edge[loop above] node[above] {\scriptsize $1:B$} (q1);
\end{tikzpicture}
        \end{center}
        Apaga o traço inicial e o reescreve em seguida, parando em $q_2$ na configuração final
        padrão; computa a identidade.
        \end{solution}
    \end{Qlist}
}

\Qitem{ \Qq{Compute a função diagonal $d(n)$ para $n = 7$. Para isto, identifique e desenhe o
        fluxograma correspondente à máquina de Turing $M_7$.}

    \begin{solution}
    $M_7$ tem número $1470 = 2\cdot 3\cdot 5\cdot 7^2$, com
    $(a_1,a_2,a_3,a_4)=(1,1,1,2)$: $q_1S_0S_0q_1$ e $q_1S_1S_0q_2$.
    \begin{center}
    \begin{tikzpicture}[>=Stealth,node distance=1cm,every node/.style={}]
  \node[tmstate] (q1) at (0.00,0.00) {$q_{1}$};
  \node[tmstate,double] (q2) at (2.80,0.00) {$q_{2}$};
  \draw[->] (q1) edge[loop above] node[above] {\scriptsize $0:B$} (q1);
  \draw[->] (q1) edge node[above,midway] {\scriptsize $1:B$} (q2);
\end{tikzpicture}
    \end{center}
    Para o argumento $7$ a fita inicial é $1\,1\,1\,1\,1\,1\,1$, com o cabeçote sobre o primeiro
    símbolo em $q_1$. A única transição aplicável, $q_1 1 \longrightarrow q_2 0$ (lê $1$, escreve o
    branco $0$ e passa a $q_2$), produz
    \[
      1_1\,1\,1\,1\,1\,1\,1 \;\longrightarrow\; 0_2\,1\,1\,1\,1\,1\,1,
    \]
    Em $q_2$ não há transição para branco, então a máquina para sobre um branco, fora da configuração
    final padrão. Logo, $f_7(7)$ é indefinida e $d(7) = 1$.
    \end{solution}
}

\section*{Parte C - Ábacos}

\begin{mdframed}
\noindent\textbf{Convenções:}\par\medskip
Os argumentos iniciam nos registradores $R_1$ e $R_2$, o resultado fica em $R_0$ e os
demais registradores começam em $0$. Um nó $r+$ incrementa $R_r$; um nó $r-$ decrementa
$R_r$ se $[r]\neq 0$, ou segue a seta $e$ se $[r]=0$.
\end{mdframed}

\Qitem{ \Qq{Especifique um ábaco para computar a função diferença `$-$' definida como segue:
        $-(x, y) = x - y$, se $y < x$, e $-(x, y) = 0$, caso contrário.}

    \begin{solution}
    \begin{center}
    \begin{tikzpicture}[>=Stealth]
        \coordinate (A) at (0, 1.4);
        \node[abnode,
            % label={[font=\tiny]135:{(1)}}
        ] (n1) at (0, 0) {$1-$};
        \node[abnode,
            % label={[font=\tiny]135:{(2)}}
        ] (n2) at (0, -2.2) {$0+$};
        \node[abnode,
            % label={[font=\tiny]135:{(3)}}
        ] (n3) at (3, 0) {$2-$};
        \node[abnode,
            % label={[font=\tiny]135:{(4)}}
        ] (n4) at (3, -2.2) {$0-$};

        \draw[->] (A) -- (n1);
        \draw[->] (n1) edge[bend right=25] (n2);
        \draw[->] (n2) edge[bend right=25] (n1);
        \draw[->] (n1) -- node[above] {\scriptsize $e$} (n3);
        \draw[->] (n3) edge[bend right=25] (n4);
        \draw[->] (n4) edge[bend right=25] (n3);
        \draw[->] (n3) -- node[above] {\scriptsize $e$} ++(2.2, 0) node[right] {\scriptsize pare};
        \draw[->] (n4) -- node[above] {\scriptsize $e$} ++(2.2, 0) node[right] {\scriptsize pare};
    \end{tikzpicture}
    \end{center}
    % O laço $(1)$--$(2)$ copia $x$ para $R_0$, e $(3)$--$(4)$ decrementa $R_2$ e $R_0$ em paralelo:
    % se $R_2$ zera primeiro, $(3)$ para com $[0]=x-y$; se $R_0$ zera primeiro ($y\geq x$), $(4)$ para
    % com $[0]=0$.
    \end{solution}
}

\Qitem{ \Qq{A função sinal $\mathrm{sn}$ é definida por $\mathrm{sn}(x) = 1$ se $x > 0$, e $= 0$ 
        caso contrário. Apresente uma prova direta de que $\mathrm{sn}$ é computável por ábaco, 
        especificando um ábaco que a compute. }

    \begin{solution}
    \begin{center}
    \begin{tikzpicture}[>=Stealth]
        \coordinate (A) at (0, 1.4);
        \node[abnode,
            % label={[font=\tiny]135:{(1)}}
        ] (n1) at (0, 0) {$1-$};
        \node[abnode,
            % label={[font=\tiny]135:{(2)}}
        ] (n2) at (0, -2.2) {$0+$};

        \draw[->] (A) -- (n1);
        \draw[->] (n1) -- (n2);
        \draw[->] (n1) -- node[above] {\scriptsize $e$} ++(2.2, 0) node[right] {\scriptsize pare};
        \draw[->] (n2) -- ++(0, -1.3) node[below] {\scriptsize pare};
    \end{tikzpicture}
    \end{center}
    % Se $x=0$, $[1]=0$ e $(1)$ segue pela seta $e$, parando com $[0]=0=\mathrm{sn}(0)$. Se $x>0$,
    % $(1)$ decrementa $R_1$ e $(2)$ incrementa $R_0$, parando com $[0]=1=\mathrm{sn}(x)$.
    \end{solution}
}

\Qitem{ \Qq{Mostre, especificando um ábaco apropriado, que a função $f$ definida por $f(x, y) = 1$,
        se $x < y$, e $f(x, y) = 0$, caso contrário, é computável por ábaco. }

    \begin{solution}
    \begin{center}
    \begin{tikzpicture}[>=Stealth]
        \coordinate (A) at (0, 1.4);
        \node[abnode,
            % label={[font=\tiny]135:{(1)}}
        ] (n1) at (0, 0) {$1-$};
        \node[abnode,
            % label={[font=\tiny]135:{(2)}}
        ] (n2) at (0, -2.2) {$2-$};
        \node[abnode,
            % label={[font=\tiny]135:{(3)}}
        ] (n4) at (3, 0) {$2-$};
        \node[abnode,
            % label={[font=\tiny]135:{(4)}}
        ] (n5) at (3, -2.2) {$0+$};

        \draw[->] (A) -- (n1);
        \draw[->] (n1) edge[bend right=25] (n2);
        \draw[->] (n2) edge[bend right=25] (n1);
        \draw[->] (n1) -- node[above] {\scriptsize $e$} (n4);
        \draw[->] (n2) -- node[above] {\scriptsize $e$} ++(-2.2, 0) node[left] {\scriptsize pare};
        \draw[->] (n4) -- (n5);
        \draw[->] (n4) -- node[above] {\scriptsize $e$} ++(2.2, 0) node[right] {\scriptsize pare};
        \draw[->] (n5) -- ++(0, -1.3) node[below] {\scriptsize pare};
    \end{tikzpicture}
    \end{center}
    % O laço $(1)$--$(2)$ decrementa $R_1$ e $R_2$ em paralelo: se $R_2$ zera primeiro ($x>y$), $(2)$
    % para com $[0]=0$; se $R_1$ zera primeiro, $(3)$ verifica se $[2]=y-x>0$, incrementando $R_0$ em
    % $(4)$ quando $x<y$ e parando com $[0]=0$ quando $x=y$.
    \end{solution}
}

\end{document}
